\chapter{Introducción}

\section{Objetivo de la práctica}

El objetivo principal de esta práctica es utilizar \textit{OpenSSL} para realizar 
operaciones criptográficas relativas al cifrado asimétrico, incluyendo:

\begin{itemize}[label=\textbullet]
    \item Generación y gestión de claves asimétricas (RSA, DSA)
    \item Exportación y conversión de formatos de claves (PEM, DER)
    \item Firma digital y verificación de resúmenes
    \item Derivación de secretos compartidos con Diffie-Hellman (DH) y curvas elípticas (X25519)
    \item Intercambio seguro de información combinando cifrado simétrico, asimétrico y firma digital
\end{itemize}

\section{Conceptos fundamentales}

\subsection{Criptografía asimétrica}

A diferencia del cifrado simétrico (donde emisor y receptor comparten la misma clave), la criptografía asimétrica utiliza un \textbf{par de claves} matemáticamente relacionadas:

\begin{itemize}[label=\textbullet]
    \item \textbf{Clave pública:} Puede distribuirse libremente. Se usa para cifrar y verificar firmas.
    \item \textbf{Clave privada:} Debe mantenerse en secreto. Se usa para descifrar y firmar.
\end{itemize}

La relación entre ambas claves cumple las siguientes propiedades:

\begin{equation}
D(K_{priv}, E(K_{pub}, M)) = M \quad \text{(cifrado)}
\end{equation}
\begin{equation}
V(K_{pub}, S(K_{priv}, H(M))) = \text{true} \quad \text{(firma)}
\end{equation}

\subsection{Algoritmos utilizados}

\begin{table}[H]
    \centering
    \caption{Algoritmos asimétricos utilizados en esta práctica}
    \begin{tabularx}{\textwidth}{l|p{3cm}|p{6cm}}
        \toprule
        \textbf{Algoritmo} & \textbf{Función} & \textbf{Descripción} \\
        \midrule
        RSA & Cifrado y firma & Basado en la factorización de números primos grandes \\
        DSA & Solo firma & Estándar de firma digital del NIST \\
        DH & Intercambio de claves & Protocolo de negociación de secretos compartidos \\
        X25519 & Intercambio de claves & DH con curva elíptica Curve25519 (usado por WhatsApp) \\
        \bottomrule
    \end{tabularx}
\end{table}

\subsection{Formatos de claves}

\begin{table}[H]
    \centering
    \caption{Formatos de almacenamiento de claves}
    \begin{tabularx}{\textwidth}{l|l|p{7cm}}
        \toprule
        \textbf{Formato} & \textbf{Tipo} & \textbf{Características} \\
        \midrule
        PEM & Texto (Base64) & Legible, puede incluir cifrado con contraseña \\
        DER & Binario & Más compacto, \textbf{no soporta cifrado con contraseña} \\
        \bottomrule
    \end{tabularx}
\end{table}

\subsection{Esquema híbrido de cifrado}

En la práctica real, el cifrado asimétrico no se utiliza directamente sobre documentos grandes debido a su lentitud. En su lugar, se emplea un \textbf{esquema híbrido}:

\begin{enumerate}
    \item Se cifra el documento con un algoritmo simétrico (rápido), por ejemplo AES-256-CBC.
    \item Se cifra la clave simétrica utilizada con la clave pública del destinatario (RSA).
    \item Se firma un resumen del documento original con la clave privada del emisor.
    \item Se envían: documento cifrado + clave cifrada + firma digital.
\end{enumerate}

Con este esquema se logra simultáneamente:
\begin{itemize}[label=\textbullet]
    \item \textbf{Confidencialidad:} Solo el destinatario puede descifrar la clave simétrica.
    \item \textbf{Integridad:} El resumen garantiza que el documento no ha sido alterado.
    \item \textbf{Autenticación:} La firma demuestra que el documento proviene del emisor legítimo.
\end{itemize}

\section{Herramientas utilizadas}

\begin{itemize}[label=\textbullet]
    \item \textbf{OpenSSL 3.x:} Herramienta principal de criptografía
    \item Comandos principales: \texttt{genpkey}, \texttt{pkey}, \texttt{pkeyutl}, \texttt{enc}, \texttt{dgst}
\end{itemize}
